\chapter{Supplementary Question 2 on Energy Gentrification in Western Sydney}
\label{ch:supp:gentrification}

\section{The question}
\label{sec:supp:gen:q}

\begin{quote}
\itshape
You discussed the concept of ``energy gentrification'' -- where concentrated data centre development drives up energy costs for local consumers -- citing examples from the United States and Ireland. Are you aware of any evidence or analysis of whether this effect is occurring or likely to occur in NSW, particularly in Western Sydney where data centre clustering is most concentrated?
\end{quote}

Protocol is not aware of any published, New South Wales specific quantitative study that calculates an energy gentrification effect, being a measurable increase in residential energy costs attributable to data centre concentration, for the State or for Western Sydney. This is stated plainly at the outset. What follows is the evidence that does exist, which establishes the concept, documents the effect in comparable jurisdictions, identifies the causal mechanism, and shows that the conditions which produced the effect elsewhere are present in Western Sydney. The conclusion drawn is bounded accordingly. The effect is a material, foreseeable and preventable risk rather than a demonstrated fact, and it is the preventability that makes early action worthwhile.

\section{The concept and its scholarly basis}
\label{sec:supp:gen:concept}

Energy gentrification belongs to a wider scholarship on digital extractivism that treats the data centre as an extractive industry rather than a neutral utility. The framing is direct, holding that \enquote{AI is a technology of extraction: from the energy and minerals needed to build and sustain its infrastructure, to the exploited workers behind \enquote{automated} services, to the data AI collects from us} \autocite{crawfordAtlasAIPower2021}. Applied to electricity, this becomes the proposition that \enquote{The global data center industry relies on what this article defines as \enquote{climate extraction}} \autocite{brodieClimateExtractionSupply2020}. The Irish case studies are the most developed, and they show how siting decisions are presented as benign while their fiscal and infrastructural logic is obscured, given that \enquote{Ireland has been advertised to and by data center developers because of its \enquote{cool} climate while downplaying the importance of its low corporate tax rate} \autocite{brodieClimateExtractionSupply2020}. The cost of that arrangement is socialised, because the new infrastructures are \enquote{entangled with non-renewable energy systems and tax evasive capital, and built across existing communities and environments} \autocite{bresnihanNewExtractiveFrontiers2021}. The governance gap is the recurring theme, captured in the observation from the Nordic case that \enquote{no one is steering this development, or closely looking at the impacts that it may have on remote communities in the Arctic and Nordic region} \autocite{sovacoolWholeSystemsEnergy2022}.

\section{International evidence}
\label{sec:supp:gen:intl}

Ireland supplies the clearest quantified case. Data centre demand there has risen from a small share to a dominant one, since \enquote{Data centre usage was only 5\% in 2015} but \enquote{Data centres consumed 22\% of Irish electricity in 2024} \autocite{irishTimesDataCentreBills2026}. The consumer cost has been estimated directly, with analysis finding \enquote{Cumulatively for all Irish households: \texteuro{}715 million in price effects} and that \enquote{An estimated cumulative average of \texteuro{}360 has been added to household electricity bills over recent years} \autocite{irishTimesDataCentreBills2026}. That cost coexists with substantial investment, as \enquote{Data centre investors have invested \texteuro{}18 billion into the Irish economy}, which is the trade-off the New South Wales inquiry must weigh rather than assume away \autocite{irishTimesDataCentreBills2026}.

Virginia, the largest data centre market in the world, supplies the second case. A legislative review found that \enquote{unconstrained demand for power in Virginia would double within the next 10 years, with the data center industry being the main driver}, and that the cost falls broadly, since \enquote{data centers' increased energy demand will likely increase system costs for all customers, including non-data center customers, for several reasons} \autocite{jlarcVirginiaDataCentres2024}. The counterfactual matters, because \enquote{The state's energy demand was essentially flat from 2006 to 2020 because, even though population increased, it was offset by energy efficiency improvements} \autocite{jlarcVirginiaDataCentres2024}, which establishes that the new load, not background growth, is the driver. Australian analysts have already drawn the comparison, noting that \enquote{In Virginia, USA, unconstrained demand from data centres is projected to add almost \$40 US to monthly residential electricity bills} \autocite{ecaConsumerImpactsDataCentres2025}.

\section{The causal mechanism}
\label{sec:supp:gen:mechanism}

The mechanism by which concentrated load can raise consumer bills is well characterised in the engineering and market literature, and it operates through three channels. The first is system scale, because a hyperscale facility is large relative to the grid that hosts it, and \enquote{With hyperscale datacenters exceeding 100 MW individually, and in some grids exceeding 15\% of power load, DC adaptation is large enough to harm power grid dynamics, increasing carbon emissions, power prices, or reduce grid reliability} \autocite{linAdaptingDatacenterCapacity2023}. The second is price formation, since the addition of large, continuous load affects \enquote{electricity markets and pricing, economic dispatch and reserve scheduling, and infrastructure planning and coordination} \autocite{shengPowerAIData2026}, against a backdrop in which \enquote{The rapid growth of artificial intelligence (AI) is driving an unprecedented increase in the electricity demand of AI data centers, raising emerging challenges for electric power grids} \autocite{chenElectricityDemandGrid2025}. The third is cost recovery, because the network augmentation and system services that new load requires are recovered from consumers, and in Australia \enquote{electricity network costs make up almost half of household electricity bills} \autocite{ecaDataCentresEnergyLandscape2025}.

Whether these channels raise residential bills depends on how cost is allocated and on whether supply responds, so the mechanism is contingent rather than deterministic. The Australian market already applies a beneficiary-pays principle to some system services, since \enquote{the costs to provide FCAS are recovered from those who create the need, such as a large electricity user} \autocite{ecaDataCentresEnergyLandscape2025}. The price pressure is also conditional on the generation response, because the load is largely inflexible and facilities frequently draw carbon-intensive supply, since \enquote{datacenters often consume carbon-intensive energy from the grid when carbon-free energy is scarce} \autocite{acunCarbonExplorerHolistic2023}. The international price signal is nonetheless striking, with reporting that \enquote{wholesale electricity costs are as much as 267\% higher today in areas near data centres than five years ago} and that, in Ireland, \enquote{data centres have been responsible for 88\% of increased electricity demand from 2015 to 2024} \autocite{ecaDataCentresEnergyLandscape2025}.

\section{The Western Sydney pipeline and grid outlook}
\label{sec:supp:gen:wsydney}

The conditions that produced the effect overseas are present in Western Sydney in concentrated form. Clustering is the first condition. The State has approved what \enquote{The biggest data centre in the Southern Hemisphere has been given the green light by the Minns Labor Government}, a facility of which \enquote{The facility is at Marsden Park in Sydney's north-west approximately 36 kilometres from the centre of Sydney} \autocite{nswMarsdenParkDataCentre2025}. Capacity concentration at the State level is acute, with New South Wales and Victoria together hosting \enquote{around 80\% of the country's data centre capacity} \autocite{usscPoweringCloud2025}. Demand growth is the second condition. The New South Wales Government's own consultation paper relays that \enquote{AEMO forecasts that energy consumption from data centres will double from 5\% of NSW's grid-supplied energy in 2026 to 11\% by 2030} \autocite{inswDataCentreConsultation2026}, a trajectory corroborated by the finding that data centres \enquote{Data centres currently account for 4\% of the state's grid-supplied electricity in NSW, expected to rise to 11\% by 2030} \autocite{usscPoweringCloud2025}. The committed pipeline alone is material, since \enquote{Data centres made up around 3 TWh of load in 2024, with the ESOO Central scenario forecasting around 5 TWh by 2033-34 from existing and committed projects alone} \autocite{ecaConsumerImpactsDataCentres2025}.

The network response is the third condition, and it is already being planned at scale. Transmission planning records that \enquote{Network Planning progressed studies to investigate capacity to support data centre load growth in Western Sydney}, and that \enquote{Studies to-date indicate additional capacity of 2.8 GW to supply load can be enabled in western Sydney, if the following conditions are met} \autocite{transgridWesternSydneyAugmentation2024}. Those conditions are extensive and are explicitly to be funded by the connecting parties, requiring that \enquote{Data centres connecting to the Transmission network fund the following augmentation projects: double-circuiting of line 14 between Kemps Creek and Sydney West with high-capacity conductors, approx. 280 MVAr dynamic voltage support, an additional 1,500 MVA 500/330 kV transformer at the Kemps Creek substation and data centre connection capacitor banks totalling 800 MVAr} \autocite{transgridWesternSydneyAugmentation2024}, and further that \enquote{Committed BESS projects along Sydney's supply corridors totalling 2.6 GW proceed as planned} \autocite{transgridWesternSydneyAugmentation2024}. The coordination risk is that these elements proceed out of step, because \enquote{Transmission investment decisions can lag data centre development timelines, and renewable projects may be delayed by network and supply challenges} \autocite{usscPoweringCloud2025}.

\section{The distribution of cost}
\label{sec:supp:gen:distribution}

The decisive variable is who pays for that augmentation. Where connection costs are reflective and borne by the connecting party, the socialisation that drives gentrification is contained; where they are not, residential consumers absorb the difference through network charges that already \enquote{make up almost half of household electricity bills} \autocite{ecaDataCentresEnergyLandscape2025}. The New South Wales Government has articulated the correct principle, that \enquote{Data centre developers and operators need to fund their infrastructure requirements, in addition to what is already planned and funded, so as to not increase prices for households}, and has proposed that \enquote{Cost recovery regimes should be reviewed at the respective national and state levels to ensure upgrades required for data centre connections are paid for by data centres} \autocite{inswDataCentreConsultation2026}. That principle is currently a stated intention rather than a binding rule, and the evidence that load growth is already entering regulated revenue proposals is concrete, since \enquote{Jemena's revenue proposal for the 2026-2031 regulatory period outlined a significant increase in non-residential electricity consumption from 2024 levels due primarily to data centre load growth} \autocite{ecaConsumerImpactsDataCentres2025}. The wholesale dimension is equally contingent, with Australian modelling indicating that unmatched data centre growth could raise New South Wales wholesale prices by around 26 per cent by 2035, an impact that falls to around 3 per cent where matching renewable generation and storage proceed \autocite{cefcBaringaDataCentres2025}.

\section{Why no New South Wales finding yet exists}
\label{sec:supp:gen:gap}

The absence of a completed New South Wales study reflects the recency of the load, not the absence of risk. The relevant research is being commissioned now, on the explicit premise that, \enquote{Based on international evidence, data centres could add hundreds of dollars annually to Australian household energy bills through a range of factors} \autocite{ecaConsumerImpactsDataCentres2025}. Until that work reports, the honest position is that the precursors are present and internationally evidenced, that the magnitude in New South Wales is not yet quantified, and that the outcome remains contingent on the cost-allocation and renewable-matching decisions the Committee and the Government are presently considering. Inferring a New South Wales figure from Irish or Virginian numbers would overstate the evidence, and this submission does not do so.

\section{Recommendation}
\label{sec:supp:gen:rec}

Protocol recommends a posture of monitoring with pre-emptive safeguards rather than retrospective remediation of the kind now under way in Ireland and Virginia. First, the consumer-bill research being scoped by Energy Consumers Australia should be commissioned with a Western Sydney regional breakdown and maintained as a standing dataset, so that the question can be answered with local evidence rather than overseas analogues \autocite{ecaConsumerImpactsDataCentres2025}. Second, the developer-funded principle already stated by the Government should be given binding legislative force, converting the intention that data centres \enquote{fund their infrastructure requirements} into an enforceable cost-recovery rule \autocite{inswDataCentreConsultation2026}. Third, consent should carry minimum safeguards consistent with the Climate Council's submission, which recommends that proponents \enquote{Support demonstrably additional renewable generation and firming} and \enquote{Disclose and report consistent information on electricity use, water use, backup generation and emissions} \autocite{climateCouncilNSWDataCentres2026}. The renewable-matching requirement is the single most effective lever, because it is the variable that the modelling shows separates a 26 per cent wholesale impact from a 3 per cent one \autocite{cefcBaringaDataCentres2025}, and because in its absence new load is liable to lean on existing or fossil supply built for other consumers \autocite{greenpeaceEnergyVampires2026}. Finally, approvals, transmission planning and renewable zone development should be coordinated rather than run as \enquote{separate processes across federal and state jurisdictions} \autocite{usscPoweringCloud2025}.
